\documentclass[sigplan,10pt,anonymous,review]{acmart}

\settopmatter{printacmref=false, printccs=false, printfolios=true}
\setcopyright{none}
\renewcommand\footnotetextcopyrightpermission[1]{}
\pagestyle{plain}
\acmConference[ATC '26]{ACM SIGOPS Annual Technical Conference}{November 16--18, 2026}{Hong Kong}
\acmYear{2026}
\acmISBN{}
\acmDOI{}

\usepackage{booktabs}
\usepackage{xspace}

\newcommand{\sysname}{\textsc{AgentScope}\xspace}

\begin{document}

\title[Extended Abstract: Telemetry-Driven Runtime Control for LLM Agents]{Extended Abstract \\
Telemetry-Driven Runtime Control for Heterogeneous LLM Agent Systems}

\author{Anonymous Authors}
\affiliation{%
  \institution{Submission under double-blind review}
  \country{}
}
\renewcommand{\shortauthors}{Anonymous Authors}

\begin{abstract}
LLM agent systems fail in ways that production tracing stacks cannot
see. \sysname{} is an OpenTelemetry SDK that closes the gap with (i) a
nine-kind agent-orchestration span vocabulary extending GenAI semantic
conventions, (ii) seven framework adapters that solve the cross-process
correlation problem, and (iii) a telemetry-driven \emph{circuit breaker}
that turns the same span stream into runtime control. Measured on
commodity hardware: 11.7\,\textmu{}s p50 per span, 19{,}071 spans/s
under 100-thread concurrency, +2.16\,\textmu{}s p50 breaker overhead
(95\% CI [+2.06, +2.25] across 5 independent runs),
and 0.612 (vs 0.429 baseline; 1.000 on a conformance-complete adapter)
fault-detection rate over a 3{,}780-row controlled benchmark spanning
7 frameworks $\times$ 6 LLMs $\times$ 14 fault classes.
\end{abstract}

\maketitle

\section*{Problem}

Production LLM agent systems --- multi-step planners that compose
reasoning, tool calls, memory operations, and inter-agent delegation
--- have entered the critical path of software engineering, customer
support, knowledge work, and incident response. The tracing stack that
was built for microservices does not see what they do. An OTel
\texttt{INTERNAL} span carries no semantic information distinguishing a
guardrail rejection from a reasoning step or a delegation handoff.
The in-stabilization OpenTelemetry GenAI semantic conventions name LLM
invocations and tool executions but stop short of the orchestration
phases that surround them. Closed-source platforms (LangSmith,
Langfuse) and adjacent OSS schemas (OpenInference, OpenLLMetry)
capture richer attributes per vendor but ship vocabularies that do not
interoperate.

The operational consequence is concrete: in a prior controlled study,
84 of 112 SWE-bench Lite runs against a production agent exhausted
their iteration budget with \emph{no} span-level evidence of why
they failed --- no record of which planning step diverged, which tool
call repeated, or which delegation cycle never closed.

This is not a documentation problem; it is a systems problem with
three dimensions: (D1) heterogeneous instrumentation surfaces across
seven popular frameworks each with a different hook model;
(D2) span correlation across asyncio / multi-process / inter-agent
boundaries; (D3) passive tracing cannot enforce runtime policies
(cost, loop, delegation cycle) without a control surface wired into
the span stream.

\section*{Approach}

\sysname{} is an OTel-native SDK addressing D1--D3 with measured
trade-offs.

\textbf{Span vocabulary.} Nine agent-specific kinds (AGENT, PLANNING,
REASONING, DELEGATION, GUARD\_RAIL, MEMORY, RETRIEVAL, TOOL\_CALL,
LLM\_CALL) implemented as OTel attributes (not new SpanKind enum
values, which would break OTLP). The taxonomy is derived from a
coverage analysis: it is the minimum closed set under which all 14
fault classes in our benchmark have a span-level signal.

\textbf{Adapter layer.} Seven framework adapters use five qualitatively
different binding strategies: callback (LangChain), lifecycle hook
(CrewAI, AutoGen), span handler (LlamaIndex), monkey-patch (Anthropic
SDK, OpenAI SDK), and manual context manager (custom). All adapters
import their framework lazily at \texttt{instrument()} time --- a hard
requirement for a single observability library that ships in agent
containers with diverse dependency closures.

\textbf{Cross-boundary correlation.} Three boundaries break OTel's
default propagation: asyncio task creation, process-pool execution,
and inter-agent message routing. We solve each with a thin wrapper
that serializes the W3C trace-context carrier and re-installs it on
the receiver side.

\textbf{Telemetry-driven circuit breaker.} An OTel
\texttt{SpanProcessor} that subscribes to the same stream as the
exporter, maintains per-trace state (cost accumulator, reasoning-step
counter, delegation graph, token-history), and fires configurable
actions (log / callback / raise) when a policy matches. The breaker
ships four default policies (reasoning loop, cost explosion, context
overflow, delegation cycle) and accepts user-defined policies through
a small Python DSL. This mechanism is impossible without the
nine-kind vocabulary: a reasoning-loop detector cannot distinguish
reasoning from delegation without per-kind attribution.

\section*{Evaluation}

We report four measurement axes. Hardware: Apple M4 Pro, Python 3.12,
opentelemetry-api/sdk 1.27.x. All raw data and reproduction commands
ship with the artifact.

\textbf{Per-span overhead.} 10{,}000 iterations per kind (90{,}000
total): aggregate p50 11.7\,\textmu{}s, p95 13.2\,\textmu{}s, p99
28.2\,\textmu{}s, mean 12.8\,\textmu{}s. Throughput 58{,}771--66{,}396
spans/s depending on kind. Memory per span 3{,}103--3{,}191 bytes.

\textbf{Scalability.} 100-thread concurrent: 19{,}071 spans/s, p50
44.0\,\textmu{}s. Long-running 1{,}200-span single trace: 7.6\%
latency degradation between first and last decile. Disk-backed JSON
exporter: 1{,}232\% blocking overhead vs in-memory baseline (we surface
this rather than hide it; the mitigation is the
\texttt{BatchSpanProcessor}, which holds 5{,}335 spans/s end-to-end on
the long trace).

\textbf{Circuit breaker activation.} Four-policy breaker installed on
top of the per-span pipeline: +2.16\,\textmu{}s p50 (+18.9\%, 95\% CI
[+2.06, +2.25] across 5 independent runs);
+3.31\,\textmu{}s p99 (+22.9\%); sub-millisecond absolute on operations
whose wall-clock baseline is 100\,ms--10\,s.

\textbf{Fault detection.} Controlled fault-injection benchmark over
7 frameworks $\times$ 6 LLMs $\times$ 14 fault classes (3{,}780 runs).
Aggregate FDR by condition: no telemetry 0.000; vanilla OTel,
OTel GenAI semconv, OpenInference all tie at 0.429
(95\% Wilson CI [0.389, 0.469]); \sysname{} METADATA\_ONLY and FULL
both 0.612 (CI [0.572, 0.651]); \sysname{} on a conformance-complete
adapter 1.000 (CI [0.957, 1.000]). Difference between \sysname{} and
production baselines significant at $p<0.001$. The per-fault matrix
shows that the three baselines catch the same 6 bottom-up classes
(tool failure, timeout, cost explosion, etc.) and miss the same 8
orchestration classes (circular delegation, reasoning loop,
guardrail bypass, etc.); the 18.3-point absolute lift is the value of
the agent-specific schema, and the gap between 0.612 and 1.000 is a
per-app conformance issue, not a metamodel-expressivity issue.

\textbf{Real-LLM end-to-end.} 45-run multi-agent topology study
(hierarchical / parallel / sequential $\times$ 3 model back-ends
$\times$ 5 seeds): 570 spans, zero orphan spans across boundaries
(Clopper--Pearson 95\% upper bound on true orphan rate: 0.52\%), span-
level cost attribution matches model provider billing within 1\%.

\section*{Contribution}

\sysname{} is the first observability stack we know of to (a) treat
agent-orchestration spans as enforceable runtime state rather than
only diagnostic exhaust, and (b) report end-to-end overhead,
scalability, and fault-detection-rate numbers for the GenAI tracing
pipeline. The systems trade-offs (sync vs batched exporter, schema
richness vs per-span memory, adapter fragility vs instrumentation
coverage, breaker action choice) are reported in full. The full paper
(under separate cover) provides per-class measurement, per-framework
breakdown, threats-to-validity analysis, and reproducibility commands.

\end{document}
